\documentclass[11pt]{report}
% T0_preamble_standalone_A4_En.tex — LOKALER PLATZHALTER
% Im Repo durch die offizielle Korpus-Präambel ersetzen; die
% \input-Zeilen der Wrapper bleiben unverändert gültig.
\usepackage{fontspec}
\usepackage[english]{babel}
\usepackage{csquotes}
\setmainfont{Inter}[BoldFont=Inter-Bold,ItalicFont=Inter-Italic,BoldItalicFont=Inter-BoldItalic]
\setmonofont{JetBrains Mono}[BoldFont=JetBrainsMono-Bold,ItalicFont=JetBrainsMono-Italic]
\usepackage{unicode-math}
\setmathfont{Libertinus Math}
\usepackage{geometry}
\geometry{a4paper,top=2.4cm,bottom=2.4cm,left=2.8cm,right=2.8cm}
\usepackage{microtype,parskip,xcolor}
\usepackage{tcolorbox}
\tcbuselibrary{breakable,skins}
\usepackage{booktabs,longtable,array,amsmath,physics,siunitx}
\sisetup{locale=US}
\usepackage{hyperref}
\hypersetup{colorlinks=true,linkcolor=blue!60!black,urlcolor=green!40!black}

% T0_preamble_patches.tex — LOKALER PLATZHALTER
% Im Repo durch die offizielle Version ersetzen.


\title{Doc.~309 -- The Scale-Anchor Problem:\\
$\Lambda$CDM and FFGFT Compared}
\author{Johann Pascher}
\date{July 2026}

\begin{document}
\maketitle
\tableofcontents

% ============================================================
\chapter*{Doc.~309 -- The Scale-Anchor Problem}
\addcontentsline{toc}{chapter}{Doc.~309 -- The Scale-Anchor Problem}
% ============================================================

% ============================================================
\section*{What this is about}
% ============================================================

Both cosmological theories — $\Lambda$CDM and FFGFT — must introduce
somewhere a number that couples the cosmic scale to the SI system.
Neither derives this anchor from first principles.

This is not a coincidence — it is a structural problem of every physical
theory that has dimensions. But the \textbf{kind and quality} of the
anchor differ fundamentally. And there is an asymmetry that has not been
explicitly named in any previous document: FFGFT inherits the systematics
of the expansion reading through its scale fit.

This document is not part of the A-series.
References: Doc.~190 (P20, P33), Doc.~267 (degeneracy),
Doc.~279 ($H_0$ chain), Doc.~308 (cosmic sector).

% ============================================================
\section*{Conventions}
% ============================================================

\noindent\textbf{Legend.}

\medskip
\begin{tabular}{@{}lp{11cm}}
\textbf{[K]}      & \emph{Core derivation} — derived mathematically or
                    numerically from the foundational posits; reproducible. \\[4pt]
\textbf{[S]}      & \emph{Structural claim} — justified but not yet fully
                    proved. \\[4pt]
\textbf{[POSIT]}  & \emph{Posit} — declared assumption, not further derived. \\
\end{tabular}

\medskip
\noindent\textbf{Central quantities.}

\medskip
\begin{tabular}{@{}lp{10.5cm}}
$\xi = 4/30000$
  & Only fundamental parameter of FFGFT [POSIT, P33].
    Form ($4/3$, $3$, $2^2$) forward-derived [K];
    magnitude ($5^4$) intuitive/empirical [POSIT]. \\[4pt]
$\bar\lambda_e = \hbar/(m_e c)$
  & Reduced Compton wavelength of the electron;
    $= 3.8616\times10^{-13}$~m. \\[4pt]
$H_0$
  & Cosmological parameter; in FFGFT:
    $H_0 = (\pi/2)\,c\,\xi^{10}/\bar\lambda_e = 66.82$~km/s/Mpc
    [K], exponent~10 [POSIT, P20]. \\[4pt]
$R_H = c/H_0$
  & Geometric maximal scale of the compact geometry —
    not an expansion horizon (Doc.~182). \\[4pt]
$\nu_\text{Cs}$
  & Caesium hyperfine transition; $= 9.192\,631\,770\times10^9$~Hz
    (exact, definition); conventional SI anchor. \\[4pt]
$\nu_e = m_e c^2/h$
  & Electron frequency; $= 1.235\times10^{20}$~Hz;
    reading-independent reference frequency. \\[4pt]
$n$
  & Integer exponent in $H_0 \propto \xi^n/\bar\lambda_e$;
    $n = 10$ is a posit (P20, Doc.~190). \\[4pt]
$\Omega_\Lambda$
  & Dark-energy density in $\Lambda$CDM;
    without reading-independent referent in the static reading. \\
\end{tabular}

% ============================================================
\chapter{The SI System — a Network without a First Anchor}
% ============================================================

\section{All units traceable to one}

Since 2019 all SI base units are defined via fixed natural constants.
The seemingly ``four independent measured values'' $c$, $h$, $e$,
$k_B$ are physically not independent — they are all traceable to one
single historically conventional anchor: the frequency of the caesium
hyperfine transition $\nu_\text{Cs}$.

\begin{center}
\begin{tabular}{lll}
\toprule
Unit & Traceability & Historical origin \\
\midrule
Metre    & $c/\nu_\text{Cs}$ (numerical factor) & Earth meridian \\
Kilogram & $h\,\nu_\text{Cs}^2/c^2$ & Prototype kilogram \\
Ampere   & $e\,\nu_\text{Cs}$ & Force definition \\
Kelvin   & $h\,\nu_\text{Cs}/k_B$ & Thermometry \\
\bottomrule
\end{tabular}
\end{center}

\section{There is no first anchor}

$\nu_\text{Cs}$ was chosen in 1967 because it was technically stable
and realisable — not because it is physically distinguished. The Rb
transition, the electron frequency, or a hundred other systems would
have served equally. All present SI numerical values would then be
different — all dimensionless ratios ($m_e/m_p$ etc.) would remain
identical.

\textbf{What is truly fixed:} only dimensionless ratios such as
$m_e/m_p \approx 1/1836$. These are unit-independent.

\textbf{Consequence:} Every theory that predicts cosmological absolute
values in SI implicitly states a relation between two frequencies —
and must set some number to do so. The FFGFT statement

\[
  H_0 = \frac{\pi}{2}\,\frac{c\,\xi^{10}}{\bar\lambda_e}
\]

is equivalent to the dimensionless statement:

\[
  \frac{H_0}{\nu_e} = \frac{\pi}{2}\,\xi^{10} = 2.79\times10^{-39}
  \qquad\text{with}\quad \nu_e = \frac{m_e c^2}{h}.
\]

The SI reference is contained entirely in the choice of $\bar\lambda_e$
as the reference scale. Without that choice, $H_0/\nu_e$ is a pure
number.

% ============================================================
\chapter{$\Lambda$CDM and its Scale Anchor}
% ============================================================

\section{Six free parameters, all empirical}

$\Lambda$CDM describes cosmology with six free parameters (Planck 2018):
$H_0$, $\Omega_b h^2$, $\Omega_\text{DM} h^2$, $n_s$, $\sigma_8$,
$\tau$. All from observation — none derived from theory.

\section{$H_0$: measured circularly}

$H_0$ is determined via a four-rung distance ladder:
parallax → Cepheids → SNe~Ia → Hubble law. Each rung introduces its
own calibration assumptions. The \textbf{Hubble tension} (Planck:
$67.4$, SH0ES: $73.0$~km/s/Mpc, $5\sigma$) is a direct consequence
of this circularity.

\section{The fine-tuning problem}

$\Lambda$ corresponds to a vacuum energy density
$\rho_\Lambda \approx 5.96\times10^{-27}$~kg/m$^3$.
QFT predicts $\rho_\text{QFT} \sim 5\times10^{96}$~kg/m$^3$.
The ratio is $10^{123}$. $\Lambda$CDM has no explanation —
it sets $\Lambda$ to the observed value without theoretical origin.

% ============================================================
\chapter{FFGFT and its Scale Anchor}
% ============================================================

\section{Structure of $\xi$ [K] and magnitude [POSIT]}

$\xi = 4/30000 = 1/(3\cdot2^2\cdot5^4)$.
According to P33 (Doc.~190):

\begin{itemize}
  \item \textbf{Form} ($4/3$, factor~3, factor~$2^2$):
        forward-derived [K]
  \item \textbf{Magnitude} ($5^4 = 625$): intuitive/empirical [POSIT]
\end{itemize}

\section{The actual fit: the exponent 10}

$\xi$ itself comes from the geometry. The \textbf{exponent~10} in
$H_0 = (\pi/2)\,c\,\xi^{10}/\bar\lambda_e$ was chosen so that
$H_0 \approx 67$~km/s/Mpc. Since $\xi \sim 1.3\times10^{-4}$,
$H_0$ jumps between neighbouring exponents by a factor of $\sim$7500:

\begin{center}
\begin{tabular}{clc}
\toprule
$n$ & $H_0$ & Assessment \\
\midrule
9  & $5.0\times10^5$~km/s/Mpc & unphysical \\
\textbf{10} & $\mathbf{66.8}$~\textbf{km/s/Mpc} & plausible \\
11 & $8.9\times10^{-3}$~km/s/Mpc & unphysical \\
\bottomrule
\end{tabular}
\end{center}

There is no continuous fitting — $n=10$ is the only integer choice that
yields a physically sensible value. The justification remains: $n=10$
gives a value compatible with $\Lambda$CDM measurements (P20, Doc.~190).

\section{The inheritance: FFGFT inherits $\Lambda$CDM systematics [K]}

$67.4$~km/s/Mpc is a $\Lambda$CDM measurement value — determined under
the assumption of the expansion reading. Since $n=10$ was calibrated to
this value, FFGFT inherits the systematics of that reading.

The ``natural'' $H_0$ from the FFGFT redshift formula
$z = e^{\xi x/\bar\lambda_e} - 1$ would be:
\[
  H_0^\text{natural} = \frac{c\,\xi}{\bar\lambda_e}
  \approx 3\times10^{33}~\text{km/s/Mpc}
  \qquad\text{(unphysical)}
\]

To obtain $\sim$67~km/s/Mpc, $\xi^{10}$ is needed instead of $\xi^1$
— and the justification of this exponent is ultimately
$\Lambda$CDM compatibility.

This is not a weakness of FFGFT alone — it is the unavoidable
consequence of the degeneracy (Doc.~267): as long as both readings are
empirically indistinguishable, every theory calibrates its parameters
to the same data.

The remedy would be a reading-independent derivation of exponent~10
from the T$^4$ geometry forward — not yet achieved (P20, Doc.~190).

\section{Cross-anchoring: what FFGFT achieves nonetheless}

That $\xi^{10}$ reproduces not only $H_0$ but also $a_0$ (inertia
transition scale, Doc.~308) — that would not be expected from a pure
$H_0$ fit. That is the cross-anchoring.

\begin{center}
\begin{tabular}{llll}
\toprule
Sector & Quantity & $\xi$-dependence & Status \\
\midrule
Cosmology & $H_0$ & $\xi^{10}$ & [K] \\
Inertia regime & $a_0$ & $\xi^{10}$ & [K]/[S] \\
Particle physics & lepton mass ladder & $\xi$-powers & [K] \\
Quantum mechanics & $K_\text{frak} = 1-100\xi$ & $\xi$ & [K] \\
Geometry & $k^* = \ln\varphi/\xi$ & $\xi$ & [S] \\
\bottomrule
\end{tabular}
\end{center}

$\Omega_\Lambda$ in $\Lambda$CDM has \textbf{zero} such cross-anchorings.

% ============================================================
\chapter{Comparison}
% ============================================================

\begin{center}
\begin{tabular}{lll}
\toprule
Property & $\Lambda$CDM & FFGFT \\
\midrule
Free scale anchors & 6--7, empirical & 1 ($\xi$, P33) \\
Theoretical origin & none & none \\
Cross-anchoring & no & yes (5+ sectors) \\
Circularity & yes ($H_0$ ladder) & no \\
Inheritance of $\Lambda$CDM systematics & inherent & via exponent~10 \\
Fine-tuning & $10^{123}$ ($\Lambda$) & absent structurally \\
Hubble tension & $5\sigma$, open & structurally absent \\
\bottomrule
\end{tabular}
\end{center}

\medskip
FFGFT exchanges six isolated parameters for one cross-anchored —
no proof, but a reduction with direction. The inheritance of the
$\Lambda$CDM scale value via exponent~10 remains the central open
weakness.

% ============================================================
\chapter{Open Points (Doc.~190)}
% ============================================================

\begin{itemize}
  \item \textbf{P20:} Form of $H_0$ as $\xi$-ratio fixed; exponent~10
        is external calibration from $\Lambda$CDM, not derived.
  \item \textbf{P33:} Magnitude of $\xi$ (factor $5^4$) intuitive/
        empirical, no forward proof. Structural parallel to P20.
  \item \textbf{Inheritance (new, Doc.~309):} Exponent~10 calibrated
        to $H_0^{\Lambda\text{CDM}}$. Remedy: reading-independent
        derivation from T$^4$ geometry — not yet achieved.
\end{itemize}

\end{document}
